% !TeX root = ../教材题库_全部条目_转写工作稿.tex
% Source: XJ-B1-C05-EX-P185-Y0596
% 第5章 / 5.3 三角函数的图象与性质 / 教材 p.185 / PDF 第192页
% OCR未校对：公式与文字均可能有误；未完成数学TeX转写；禁止直接发给学生
% 原题及相邻原解截图（核对用）：../题目原图/XJ-B1-C05-EX-P185-Y0596-01.webp
% 本文件只在转写完成后作为题干使用；OCR 条目尚不是可靠题干。
\begin{bankitem}{5.3 例题 例题5}{第5章 / 5.3 三角函数的图象与性质 / 教材 p.185 / PDF 第192页}{XJ-B1-C05-EX-P185-Y0596}{OCR未校对：公式与文字均可能有误；未完成数学TeX转写；禁止直接发给学生}
\begin{ocrdraft}
[教材p.185 / PDF第192页]
\par
00求函数！y一t“（2，十的定义域和单调区间．
\par
解要使函数！y一t“（2，十）有意义，自变量应满足
\par
2x十一一十差丌（差），
\par
x乒一十一丌（差ez）．
\par
122
\par
于是函数的定义域是xxeR且x讵十丌，kez．
\par
由一一十差丌<2工十一<一十差丌，得
\par
5
\par
丌十一丌<x<一十
\par
，差eZ.
\par
12
\par
122
\par
因此，函数的单调递增区间是（一
\par
5
\par
丌十
\par
一丌，kGZ.
\par
12
\par
122
\end{ocrdraft}
\end{bankitem}
